\begin{textsample}{POS Dim 1 – go – Score 44.00 – go/go\_ef\_8.txt}  \label{ex:f1_pos_103}
3.2.
História A História é a ciência \textbf{que} estuda o \textbf{homem} em sociedade e suas ações no tempo e no espaço.
A ação do \textbf{homem} no tempo, como objeto de análise desta ciência, propicia a construção de \textbf{um} conhecimento histórico, metodologicamente orientado, \textbf{uma} vez \textbf{que} a relação passado-presente não se processa de forma automática, mas \textbf{exige} o conhecimento de referenciais \textbf{teóricos} capazes de atribuir sentido aos objetos históricos selecionados.
O conhecimento sobre o passado é também \textbf{um} conhecimento do presente, objetivado para o futuro, elaborado por distintos \textbf{sujeitos}.
Para a apropriação do conceito de temporalidade é fundamental compreender as relações entre anterioridade e posteridade, sucessão e simultaneidade, permanências e transformações, continuidades, descontinuidades e rupturas.
Esse movimento permite ao estudante a \textbf{percepção} das diversas temporalidades no curso da \textbf{humanidade}, a partir da sua existência e da sua história local, regional e nacional, visando compreender as diversas formas de organizações políticas, econômicas e socioculturais bem como o seu lugar no mundo.
A memória histórica é instrumento importante na busca por apreensão das ações dos agentes sociais.
Ela vai além das fontes escritas, ampliando e fortalecendo a compreensão \textbf{que} se pode ter acerca da \textbf{humanidade} e suas manifestações sociais.
Portanto, reconstituir o legado e herança \textbf{supõe} lidar com a memória \textbf{enquanto} história viva e \textbf{vivida} \textbf{que} permanece no tempo.
No quadro de referências simbólicas, a memória, no confronto de pluralidade de subjetividades, possibilita publicizar os acontecimentos \textbf{que} foram \textbf{relegados} aos esquecimentos ou aos silenciamentos.
A \textbf{contemporaneidade}, em toda a sua complexidade e multiplicidade de \textbf{atores} e práticas, tem ampliado os instrumentos para interrogar e oferecer respostas ao nosso mundo.
Disto é possível depreender o alargamento do arcabouço de fontes históricas e suas formas de analisá -las.
Ampliando, também, as possibilidades do ensino-aprendizagem histórico em sala de aula.
Numa sociedade cada vez mais pragmática e utilitarista dos saberes, marcada pela instantaneidade das trocas de informação, o conhecimento histórico tende a secundarizar -se.
Este conhecimento é indispensável para \textbf{que} crianças e jovens \textbf{vivam} em sociedade ao transformá -lo em saber \textbf{sistematizado}, possibilitando \textbf{uma} globalização das relações humanas e o mundo \textbf{que} os rodeia.
Tal necessidade é contemplada na primeira competência geral da BNCC, \textbf{que} é a utilização dos saberes e conhecimentos \textbf{historicamente} construídos “ sobre o mundo físico, social, cultural e digital para entender e explicar a realidade, continuar aprendendo e colaborar para a construção de \textbf{uma} sociedade \textbf{justa}, democrática e inclusiva ” ( BRASIL, 2017, p.09 ).
Dessa forma, qual a função da História para a formação do estudante na etapa do Ensino Fundamental?
As questões \textbf{que} nos levam a pensar \textbf{historicamente} como \textbf{um} saber necessário para a formação das crianças e jovens na instituição escolar, são originárias do tempo presente, de forma contextualizada e valorizando o protagonismo do estudante e do professor nas práticas pedagógicas.
Ao considerarmos a necessidade de explicação do mundo em \textbf{que} \textbf{vivemos}, em seus múltiplos sentidos, é fundamental \textbf{que} a relação passado/presente impulsione a dinâmica do ensino-aprendizagem.
Tal processo acontece, dialogicamente, na busca do desenvolvimento formativo do \textbf{sujeito} de direito, a partir de seu arcabouço de experiências socioculturais vivenciadas.
Este exercício possibilita ao estudante analisar e compreender os significados de diferentes objetos, lugares, circunstâncias, temporalidades, movimentos de pessoas, coisas e saberes, relacionando sua existência social, política, econômica, cultural e \textbf{identitária} com o mundo dinâmico e \textbf{globalizado} em \textbf{que} \textbf{vive}.
O Documento Curricular para Goiás ( DC-GO ), a partir da Base Nacional Comum Curricular ( BNCC ), procurou estabelecer, dentro de \textbf{uma} rede múltipla e heterogênea, \textbf{uma} relação de composição com o ato científico de historiar.
Mas, buscou manter -se \textbf{neutro} em apontar correntes filosóficas e \textbf{teóricas}, \textbf{métodos} e processos, referências e citações para o ensino-aprendizagem de História.
Apesar de entender a importância de tais indicações para o professor, \textbf{que} sempre \textbf{atenta} -se em saber quais foram as referências \textbf{que} nortearam a escrita, tanto dos textos quanto das habilidades, e as escolhas filosóficas/teóricas \textbf{norteadoras} do componente História, no DC-GO, optou -se por não destacar ou enfatizar tais dados e informações por dois motivos.
Primeiro, a BNCC optou por não indicar as correntes \textbf{teóricas} \textbf{norteadoras}, os \textbf{autores} e as referências \textbf{que} \textbf{influenciaram} a sua escrita, da mesma forma, seguindo o seu modelo, o componente curricular História no DC-GO não \textbf{explicita} tais sinalizações.
Segundo, decidiu -se a não sinalizar as referências no DC-GO, visando permitir a total liberdade e respeito ao professor \textbf{que} poderá lançar mão de suas escolhas \textbf{teóricas}, dos caminhos e processos conforme sua formação e \textbf{fundamentos} \textbf{teóricos}.
Caso fosse sinalizado qualquer referência ou filosofia, corrente \textbf{teórica} ou mesmo metodologias, poderia imprimir imposição ou normatização de tais em detrimentos de outros.
Como não foram sinalizados ou apontados nenhum, confiou -se nas mãos dos docentes a escolha e utilizações das \textbf{que} mais lhe parecer apropriado e \textbf{atual} para desenvolver as competências e habilidades com o estudante.
Atendendo a esses pressupostos e sua articulação com as competências gerais da BNCC e com as competências específicas da área de Ciências Humanas, o componente curricular História deve garantir aos estudantes o desenvolvimento de competências específicas \textbf{que} são : QUADRO 11 – COMPETÊNCIAS ESPECÍFICAS DE HISTÓRIA PARA O ENSINO FUNDAMENTAL | Nº | Competência | |---|---| | 1 | Compreender acontecimentos históricos, relações de poder e processos e mecanismos de transformação e manutenção das estruturas sociais, políticas, econômicas e culturais ao longo do tempo e em diferentes espaços para analisar, posicionar -se e intervir no mundo \textbf{contemporâneo}. | | 2 | Compreender a historicidade no tempo e no espaço, relacionando acontecimentos e processos de transformação e manutenção das estruturas sociais, políticas, econômicas e culturais, bem como problematizar os significados das lógicas de organização cronológica. | | 3 | Elaborar questionamentos, hipóteses, argumentos e proposições em relação a documentos, interpretações e contextos históricos específicos, recorrendo a diferentes linguagens e mídias, exercitando a empatia, o diálogo, a resolução de conflitos, a cooperação e o respeito. | | 4 | Identificar interpretações \textbf{que} expressem visões de diferentes \textbf{sujeitos}, culturas e povos com relação a \textbf{um} mesmo contexto histórico, e posicionar -se criticamente com base em princípios éticos, democráticos, inclusivos, sustentáveis e solidários. | | 5 | Analisar e compreender o movimento de populações e mercadorias no tempo e no espaço e seus significados históricos, levando em conta o respeito e a solidariedade com as diferentes populações. | | 6 | Compreender e problematizar os conceitos e procedimentos \textbf{norteadores} da produção historiográfica. | | 7 | Produzir, avaliar e utilizar tecnologias digitais de informação e comunicação de modo crítico, ético e responsável, compreendendo seus significados para os diferentes grupos ou estratos sociais. | Fonte : BNCC, 2017, p. 400.
As competências específicas do componente História devem ser desenvolvidas de maneira integrada com as competências da área de Ciências Humanas e com as competências gerais.
A seguir, tem -se \textbf{um} exemplo de \textbf{uma} das muitas possibilidades de articulação entre elas : Para \textbf{que} os estudantes desenvolvam todas as competências, visando a sua formação integral, as habilidades da BNCC foram desdobradas e organizadas, no DC-GO, em habilidades com diferentes graus de complexidades, com ampliação de escala e \textbf{percepção}.
Estes \textbf{desdobramentos} se deram de duas formas principais, sendo a primeira alinhada ao cuidado em apresentar a habilidade obedecendo a \textbf{uma} gradação cognitiva ; a segunda, contextualizar as habilidades para atender às especificidades goianas, regionais, às diversidades culturais, às múltiplas configurações \textbf{identitárias}, étnico-identitária, raciais, culturais, religiosas, sexuais e ainda contemplar os temas \textbf{atuais} na \textbf{contemporaneidade}.
No contexto goiano e regional, \textbf{preocupou} -se em destacar as \textbf{abordagens} dos povos indígenas, ciganos, descendentes africanos, comunidades de descendentes imigrantes de várias partes do mundo \textbf{que} se encontram radicados no estado de Goiás.
A busca não foi apenas em apresentá -los, mas intencionou -se trazer à tona as suas contribuições políticas, culturais, científicas e sociais na formação da sociedade brasileira, bem como na construção da goianidade, para \textbf{que} o estudante compreenda a \textbf{inter-relação} e a interdependência dos fatos micro e macro na construção da teia do saber histórico.
Nessa perspectiva, o ensino de História nos anos iniciais apresentado neste documento, visa o letramento histórico e busca envolver os estudantes no seu contexto, para a valorização de sua própria história, alargando progressivamente para a história nacional e do mundo.
Nesta fase é de \textbf{suma} importância valorizar os campos de experiências da Educação Infantil, principalmente “ O eu, o outro e o nós ”, para a ampliação da construção da noção de identidade, estabelecendo relação entre identidades individuais e sociais, \textbf{enquanto} agente atuante na história.
Dessa forma, podemos \textbf{dizer} \textbf{que} este processo \textbf{que} inicia nos campos de experiência na Educação Infantil contribui para a tomada de consciência da existência de \textbf{um} “ Eu ” e de \textbf{um} “ Outro ” pela criança, \textbf{que} vai sendo ampliada à medida \textbf{que} ela desenvolve a capacidade de administrar a sua \textbf{vontade} com autonomia, como parte de \textbf{uma} família, \textbf{uma} comunidade e \textbf{um} corpo social.
Sendo assim, torna -se \textbf{imprescindível} \textbf{que} o ensino de História permita \textbf{que} as crianças se compreendam a partir de suas próprias representações, da época em \textbf{que} \textbf{vivem}, inseridos num grupo e, ao mesmo tempo, resgatem a diversidade e reflitam sobre a memória \textbf{que} é transmitida.
Nos anos iniciais do Ensino Fundamental, as unidades temáticas \textbf{focaram} -se no reconhecimento do “ Eu ”, do “ Outro ” e do “ Nós ” e ampliaram -se para a noção do espaço e lugar em \textbf{que} \textbf{vive} e as dinâmicas em torno da cidade, diferenciando a vida privada e a pública, a urbana e a rural e ainda a circulação dos primeiros grupos humanos.
Após a concretização desta fase, essa análise alarga -se, com ênfase em pensar a diversidade dos povos e culturas e suas formas de organização, a noção de cidadania, os direitos e deveres e o reconhecimento da diversidade das sociedades, \textbf{que} pressupõem \textbf{uma} educação \textbf{que} estimule o convívio e o respeito entre os povos.
Ao priorizar a relação “ Eu ”, o “ Outro ” e o “ Nós ”, o DC-GO priorizou contextualizar o processo histórico, a partir da realidade do estudante, na busca da formação do ser humano global \textbf{que} saiba interagir com sua realidade com criticidade e autonomia.
O processo de \textbf{ensino-aprendizagem}, visando a investigação histórica na fase dos anos finais na BNCC e no DC-GO, está pautado em três procedimentos básicos : identificação da construção e formação da sociedade ocidental, rupturas e permanências ; a produção e o fazer histórico como \textbf{uma} ciência ; a capacidade de interpretação de diferentes versões, de forma retórica e discursiva, sobre o mesmo fenômeno histórico.
Estes três procedimentos, na BNCC, partem do geral para o particular, começando da história mundial e seus \textbf{desdobramentos} na história nacional.
O DC-GO, na busca de subsidiar a construção do estudante como ser integral, estreita ainda mais tal recorte ao destacar a história regional, goiana e local.
As temáticas abordadas no DC-GO, nos anos finais, aprofundam a partir do 6º ano, os conhecimentos sobre o saber e o fazer História e o começo da civilização humana, de forma harmônica e contínua com os aprendizados dos anos iniciais.
Desta forma, amplia -se, a partir do registro das primeiras sociedades e a construção da Antiguidade Clássica, com a necessária contraposição com outras sociedades e concepções de mundo, chegando até o medievo europeu e às formas de organização social e cultural, com destaque na África mediterrânea, Oriente Médio e nas Américas Pré-colombianas, no Brasil e em Goiás.
No 7º ano, as unidades temáticas perpassam pelas conexões entre Europa, América e África, a construção da ideia de modernidade e o novo mundo ante ao mundo medieval.
Também são tratados neste ano aspectos políticos, sociais, econômicos e culturais ocorridos, a partir do final do \textbf{século} XV, até o final do \textbf{século} XVIII.
No 8º ano, as unidades temáticas tratam sobre a África, a Ásia e Europa do \textbf{século} XIX e os movimentos como o nacionalismo, o imperialismo e as resistências a esses discursos e práticas são apresentados para compor a conformação histórica do mundo \textbf{contemporâneo} e fazem relação com os múltiplos processos \textbf{que} \textbf{desencadearam} as independências nas Américas, no Brasil seus \textbf{desdobramentos}, o Império Brasileiro, com ênfase em Goiás.
No 9º ano consolida -se a República Brasileira e sua \textbf{contextualização} mundial é o foco do ato de historicizar a \textbf{contemporaneidade}.
Para compreender as mudanças ocorridas no Brasil, após a proclamação da República, e o protagonismo de diferentes grupos e \textbf{atores} históricos neste período, faz -se necessário compreender os processos ocorridos na História Geral dos \textbf{séculos} XX e XXI, reconhecendo as especificidades e aproximações entre diversos eventos.
Destacam -se neste período os conflitos mundiais e nacionais, os regimes ditatoriais, o movimento socialista no mundo, a guerra fria e suas implicações nas configurações do mundo ocidental e oriental.
Tais conjunturas permitem ao estudante a compreensão circunstanciada das razões \textbf{que} presidiram o acirramento das identidades nos dias \textbf{atuais} e explicam a importância do debate sobre Direitos Humanos, com a ênfase nas diversidades.

% matched lemmas: abordagem, atentar, ator, atual, autor, contemporaneidade, contemporâneo, contextualização, desdobramento, desencadear, dizer, enquanto, ensino-aprendizagem, exigir, explicitar, focar, fundamento, globalizado, historicamente, homem, humanidade, identitária, imprescindível, influenciar, inter-relação, justo, método, neutro, norteador, percepção, preocupar, que, relegar, sistematizar, sujeito, sumo, supor, século, teórico, um, vivar, viver, vontade
\end{textsample}
